\section{\label{sec:1} Introduction}

% --- P1: spectral cues are monaural, idiosyncratic and relearnable ----------
The filtering of sound by the pinnae, head and upper body imposes
direction-dependent spectral cues that listeners use to judge the elevation of
a source and to tell front from back
\citep{blauert_sound_1969, middlebrooks_narrow-band_1992, baumgartner_modeling_2014}.
Unlike interaural time and level differences, which arise from comparing the
two ears, spectral cues are available to each ear on its own, and the mapping
from spectrum to direction is specific to the listener: it depends on the shape
of the ears, which changes with growth and with age
\citep{keen_growth_1988, otte_age-related_2013}. The mapping therefore has to
be acquired and kept calibrated, and adult listeners retain the ability to do
so. When the spectral cues are altered by molds inserted into the concha,
elevation localization collapses and then recovers over days to weeks of
exposure \citep{hofman_relearning_1998, van_wanrooij_relearning_2005, trapeau_fast_2016},
more quickly with sensorimotor training
\citep{parseihian_rapid_2012, carlile_relearning_2014}, and without impairing
localization with the native ears once the molds are removed
\citep{hofman_relearning_1998, trapeau_fast_2016}; see
\citet{carlile_plastic_2014} for a review. We have recently shown that the
same listeners can acquire two successive sets of new spectral cues
\citep{friedrich_adaptation_2025}.

% --- PLANNED (not written yet; one line each) --------------------------------
% --- P2: two stages at which a new mapping could be stored -----------------
How the two ears' spectral cues are combined constrains where such learning can
take place. A mold in one ear degrades elevation judgments mainly on the side of
that ear, the deficit fading toward the opposite side
\citep{hofman_binaural_2003, van_wanrooij_relearning_2005}. This pattern is
captured by models in which the spectrum at each ear is compared with a stored
representation of that ear's transfer functions, and the two monaural estimates
are then combined with a weight that favors the ear nearer the source
\citep{hofman_binaural_2003, baumgartner_modeling_2014}; the weight follows the
perceived lateral angle and is governed by interaural time and level
differences \citep{macpherson_binaural_2007}. In such a scheme a new
spectral-to-spatial mapping could be acquired at either of two stages: before
the combination, as a change to the stored representation of the perturbed
ear, or after it, in a representation of space that no longer records which ear
supplied the cue. \citet{van_wanrooij_relearning_2005} adapted listeners to a
mold in one ear over several weeks and concluded that learning acted
predominantly on the spectral-to-spatial mapping rather than on the binaural
weights: localization with the adapted cues improved even when the other ear
was molded as well. Whether the newly learned mapping belongs to the ear that
acquired it was not tested, however, because the learned cue was never
presented to the other ear.

% --- PLANNED (continued) -----------------------------------------------------
% P3  Why molds cannot answer this: binaural molds change both ears; a
%     monaural mold changes one ear, but the other ear's pinna differs, so the
%     learned cue can never be presented to it. Specificity as a locus probe:
%     eye-specific perceptual learning placed texture learning at a monocular
%     stage (Karni & Sagi 1991, PNAS). Brief: interaural generalization of
%     auditory learning in other tasks (verify before citing).
% P4  Approach: virtual auditory environment; own envelope + donor detail in
%     one ear, ITD/ILD intact; four days of monaural sensorimotor training;
%     final 2x2 Ear x Side, cell D physically identical to A with channels
%     swapped.
% P5  Predictions per cell: ear-specific -> A (and B, weighted by far-ear
%     weight), not C/D; shared spectral mapping -> D mirrors A; side-of-space
%     -> A and C. dB vs dC near the midline discriminates ear from side.
